\documentclass{handout}
\usepackage{handoutSetup}\usepackage{handoutShortcuts}

\newtheorem{Fact}{Fact}

\begin{document}

\begin{verbatimwrite}{\jobname.title}
Consumption with Constant Absolute Risk Aversion (CARA) Utility
\end{verbatimwrite}

\handoutHeader

\centerline{\Large \input \jobname.title}\vspace{.15in}

Consider the optimization problem of a consumer with a constant
absolute risk aversion instantaneous utility function $\uFunc(\CRat)=
-(1/\CARA) e^{-\CARA \CRat}$ implying $\uFunc^{\prime}(\CRat) =
e^{-\CARA \CRat}$ facing an interest rate that is constant at
$\rfree=\Rfree-1$.\footnote{A problem like this was considered in a well-known paper by \cite{caballero:jme}.}  The consumer's optimization problem is
\begin{equation}
\max_{\{\CFunc\}_{t}^{T}}~~ \Ex_{t}\left\{\sum_{s=t}^{T} \Discount^{s-t} \uFunc(\CRat_{s})\right\}            \label{eq:maxprob}
\end{equation}
subject to the constraints \opt{MarginNotes}{\marginpar{\tiny Introducing $K$ and $X$ because useful to have a term for non-income resources at beginning of period.}}{}
\begin{equation}\begin{gathered}\begin{aligned}
    B_{t+1} & =  ({M}_{t}-\CRat_{t})\Rfree
\\      {M}_{t+1} & =  B_{t+1}+Y_{t+1}
\end{aligned}\end{gathered}\end{equation}
where $Y_{t+1}$ is the consumer's idiosyncratic income, which exhibits a
random-walk deviation from an exogenously-growing trend:
\begin{equation*}\begin{gathered}\begin{aligned}
        \bar{\PLev}_{t+1} & =  \PGro \bar{\PLev}_{t}  \\
        Y_{t+1} & =  \bar{\PLev}_{t+1}+\PLev_{t+1}
\\      \PLev_{t+1} & =  \PLev_{t}+\PShk_{t+1}.
\end{aligned}\end{gathered}\end{equation*}

\noindent Bellman's equation for this problem is \opt{MarginNotes}{\marginpar{\tiny Bellman's eqn relates $\VFunc_{t}$ and $\VFunc_{t+1}$ through the controls and states.  Doesn't necessarily require writing bud constr.}}{}
\begin{equation}\begin{gathered}\begin{aligned}
        \VFunc_{t}({M}_{t},\bar{\PLev}_{t},\PLev_{t}) & =  \max_{\{\CFunc\}_{t}^{T}} ~~ \uFunc(\CRat_{t}) + \Ex_{t}[\Discount \VFunc_{t+1}({M}_{t+1},\bar{\PLev}_{t+1},{\PLev}_{t+1})]    \label{eq:vmax}.
\end{aligned}\end{gathered}\end{equation}

The first order condition (FOC) for the CARA utility problem is
\begin{equation}\begin{gathered}\begin{aligned}
        \uFunc^{\prime}(\CRat_{t}) & =  \Rfree \Discount \Ex_{t}[\VFunc_{t+1}^{M}]
\end{aligned}\end{gathered}\end{equation}
and the \handoutC{Envelope} theorem tells us that
\begin{equation}\begin{gathered}\begin{aligned}
        \VFunc^{M}_{t} & =  \Rfree \Discount \Ex_{t}[\VFunc_{t+1}^{M}].
\end{aligned}\end{gathered}\end{equation}

In the perfect foresight version of the model in
which $\PShk_{t} = 0 ~\forall~t$, the Euler equation will be
\begin{equation}\begin{gathered}\begin{aligned}
        \uFunc^{\prime}(\CRat_{t}) & =  \Rfree\Discount \uFunc^{\prime}(\CRat_{t+1}) \\
    \exp[-\CARA \CRat_{t}] & =  \Rfree \Discount \exp[-\CARA \CRat_{t+1}]
\\      1 & =  \Rfree\Discount \exp[-\CARA (\CRat_{t+1}-\CRat_{t})]
\\  \exp[\CARA (\CRat_{t+1}-\CRat_{t})] & =  \Rfree\Discount \label{eq:cgrow}
\\ \CARA (\CRat_{t+1}-\CRat_{t})  & =  \log \Rfree\Discount
\\  \CRat_{t+1} & =  \CRat_{t}+\log (\Rfree\Discount)^{1/\CARA}
.
\end{aligned}\end{gathered}\end{equation}

The $\log (\Rfree\Discount)^{1/\CARA}$ term reflects the intertemporal 
substitution factor in consumption.  Notice that intertemporal 
substitution takes the form of {\it additive} changes in the {\it 
level} of consumption in the CARA utility model, rather than 
multiplicative changes that affect the growth rate of consumption, as 
in the CRRA model.

Now suppose we are interested in the case where permanent income shocks are distributed 
normally, $\PShk_{t} \sim \mathcal{N}(0,\sigma_{\PShk}^{2})$.  Then it turns
out that the process
\begin{equation}\begin{gathered}\begin{aligned}
        \CRat_{t+1} & =  \CRat_{t} + \log (\Rfree\Discount)^{1/\CARA} + \CARA \sigma^{2}_{\PShk}/2+\PShk_{t+1} \label{eq:soln}
\end{aligned}\end{gathered}\end{equation}
satisfies the FOC under uncertainty:
\begin{equation}\begin{gathered}\begin{aligned}
    1 & =  \Rfree \Discount \Ex_{t}[\exp[-\CARA ({c}_{t+1}-c_{t})]]
\\  1 & =  \Rfree \Discount \Ex_{t}[\exp[-\CARA (\CARA \sigma_{\PShk}^{2}/2+{\PShk}_{t+1}+(1/\CARA)\log (\Rfree\Discount) +c_{t}-c_{t})]]    
\\  1 & =  \Rfree \Discount \exp[-\CARA^{2} \sigma_{\PShk}^{2}/2] \Ex_{t} \{\exp[-\CARA {\PShk}_{t+1}]\}\exp[-\CARA (1/\CARA) \log \Rfree\Discount]
\\  1 & =  \Rfree \Discount \exp[-\CARA^{2}(\sigma_{\PShk}^{2}/2)]\exp[\CARA^{2}(\sigma_{\PShk}^{2}/2)]\exp[\log (\Rfree\Discount)^{-1}]
\\  1 & =  \Rfree \Discount (\Rfree\Discount)^{-1}
\\  1 & =  1.
\end{aligned}\end{gathered}\end{equation}

Define $\kappa = \log (\Rfree\Discount)^{1/\CARA} + \CARA \sigma^{2}_{\PShk}/2$, so that
\eqref{eq:soln} becomes:
 \begin{equation}\begin{gathered}\begin{aligned}
        \CRat_{t+1} & =  \CRat_{t} + \PShk_{t+1} + \kappa.
\end{aligned}\end{gathered}\end{equation}

The expected present discounted value of consumption is \opt{MarginNotes}{\marginpar{\tiny The $\PShk_{t+n}$ terms disappear when expectations are taken.}}{}
\begin{equation}\begin{gathered}\begin{aligned}
        \PDV_{t}(\CRat) & =  \CRat_{t}+(\CRat_{t}+\PShk_{t+1}+\kappa)/\Rfree + (\CRat_{t}+\PShk_{t+1}+\kappa + \PShk_{t+2} + \kappa)/\Rfree^{2} + \ldots  \nonumber
\\      \Ex_{t}[\PDV_{t}(\CRat)] & =  \CRat_{t}+\CRat_{t}/\Rfree+\CRat_{t}/\Rfree^{2}+ \ldots + \kappa/\Rfree + 2\kappa/\Rfree^{2} + 3\kappa/\Rfree^{3} + \ldots
\\       & =  \CRat_{t}(1+\Rfree^{-1}+\Rfree^{-2}+\ldots) + \kappa \sum_{i=1}^{\infty} i/\Rfree^{i} .
\end{aligned}\end{gathered}\end{equation}

Now we need the following fact:
\begin{Fact} If $\Rfree > 1$, then $\displaystyle \sum_{i=0}^{\infty} i/\Rfree^{i} = \left(\frac{\Rfree}{(\Rfree-1)^{2}}\right)$ \label{fact:fact3}
\end{Fact}


Thus, the expectation of the infinite horizon PDV of consumption is:
\begin{equation}\begin{gathered}\begin{aligned}
        \Ex_{t}[\PDV_{t}(\CRat)] & =  \CRat_{t}\left(\frac{1}{1-1/\Rfree}\right)+\left(\frac{\kappa \Rfree}{(1-\Rfree)^{2}}\right)  .
\end{aligned}\end{gathered}\end{equation}

Given the process for income described above, we have
\begin{equation}\begin{gathered}\begin{aligned}
        \PDV_{t}(Y) & =  Y_{t}+Y_{t+1}/\Rfree+Y_{t+2}/\Rfree^{2}+\ldots  \\
         & =  \bar{\PLev}_{t}+\PLev_{t}+(\PGro\bar{\PLev}_{t}+\PLev_{t+1})/\Rfree + (\PGro^{2}\bar{\PLev}_{t}+\PLev_{t+2})/\Rfree^{2}+ \ldots
\\   & =  \bar{\PLev}_{t}\left(1+\PGro/\Rfree + (\PGro/\Rfree)^{2} + \ldots \right)+
\\   &    \PLev_{t}+(\PLev_{t}+\PShk_{t+1})/\Rfree+(\PLev_{t}+\PShk_{t+1}+\PShk_{t+2})/\Rfree^{2}+\ldots    
\\ \Ex_{t}[\PDV_{t}(Y)]   & =  \left(\frac{\bar{\PLev}_{t}}{1-\PGro/\Rfree}\right)+\PLev_{t}\sum_{s=0}^{\infty} \Rfree^{-s}
\\   & =  \left(\frac{\bar{\PLev}_{t}}{1-\PGro/\Rfree}\right)+\left(\frac{ \PLev_{t}}{1-1/\Rfree}\right)
\end{aligned}\end{gathered}\end{equation}

The IBC says
\begin{equation}\begin{gathered}\begin{aligned}
        \PDV_{t}(\CRat) & =  B_{t}+\PDV_{t}(Y),
\end{aligned}\end{gathered}\end{equation}

Because the intertemporal budget constraint must hold in every state 
of the world, the expectation of the PDV of consumption must equal current
wealth plus the expectation of the PDV of income.  Thus,
\begin{equation}\begin{gathered}\begin{aligned}
        \Ex_{t}[\PDV_{t}(\CRat)] & =  B_{t}+\Ex_{t}[\PDV_{t}(Y)  ]
\\ \CRat_{t}\left(\frac{1}{1-1/\Rfree}\right)    & =  B_{t}+ \left(\frac{\bar{\PLev}_{t}}{1-\PGro/\Rfree}\right)+\left(\frac{ \PLev_{t}}{1-1/\Rfree}\right)-\left(\frac{\kappa \Rfree}{(1-\Rfree)^{2}}\right) 
\\ \CRat_{t}        & =  \PLev_{t} + \left(\frac{\rfree}{\Rfree}\right)\left[ B_{t}+ \left(\frac{\bar{\PLev}_{t}}{1-\PGro/\Rfree}\right)-\left(\frac{\kappa \Rfree}{(1-\Rfree)^{2}}\right)  \right]
\\ & =  \PLev_{t} + \left(\frac{\rfree}{\Rfree}\right)\left[ B_{t}+ \left(\frac{\bar{\PLev}_{t}}{1-\PGro/\Rfree}\right)\right] -\rfree\left(\frac{\log (\Rfree\Discount)^{1/\CARA} + \CARA \sigma^{2}_{\PShk}/2 }{(1-\Rfree)^{2}}\right)  \nonumber
\end{aligned}\end{gathered}\end{equation}

The $\PLev_{t}$ term reflects the consumer's idiosyncratic level of 
permanent income, which has no systematic growth (or decline).  The 
next term reflects the MPC out of total `certain' wealth, human and 
nonhuman.  The final term reflects the combination of the 
intertemporal substitution motive (in the $\log (\Rfree\Discount)^{1/\CARA}$ 
term) and the precautionary motive in the $\CARA 
\sigma^{2}_{\PShk}$ term, as is evident from the fact that it 
equals zero if either there is no precautionary motive ($\CARA=0$) or 
there is no uncertainty $\sigma^{2}_{\PShk}=0$.

Note some peculiar aspects of this solution.  First, observe that, marginally, the 
consumer spends exactly the interest income on capital, $d \CRat_{t}/d 
B_{t} = \rfree/\Rfree$.  The reason this is peculiar is that the MPC out of 
capital {\it does not depend on how impatient the consumer is}.  
Impatience is reflected in the change in consumption over time, but 
not in the level of consumption except as that is affected by the 
budget constraint.

Second, notice that the effect of income uncertainty on saving is the 
same in absolute dollars regardless of the level of resources or 
permanent income.

\input handoutBibMake.tex

\end{document}

If $t+1 = T$ so that $c_{t+1} = (m_{t}-c_{t})\Rfree+\PShk_{t+1}$ then
the Euler equation becomes
\begin{equation}\begin{gathered}\begin{aligned}
        \uFunc^{\prime}(c_{t}) & =  \Rfree\Discount \Ex_{t}[\uFunc^{\prime}(c_{t+1})] \\
    \exp[-\CARA c_{t}] & =  \Rfree \Discount \Ex_{t}[\exp[-\CARA {c}_{t+1}]]
\\  1 & =  \Rfree \Discount \Ex_{t}[\exp[-\CARA ((m_{t}-c_{t})\Rfree+{\PShk}_{t+1}-c_{t})]]
\\  0 & =  \log (\Rfree\Discount) + \log \Ex_{t}[\exp[-\CARA (-c_{t}(1+\Rfree) + Rm_{t})]\exp[-\CARA {\PShk}_{t+1}]]
\\  0 & =  \log (\Rfree\Discount) + (\CARA c_{t}(1+\Rfree) - \CARA \Rfree m_{t})+ \log \Ex_{t}[\exp[-\CARA {\PShk}_{t+1}]] \nonumber
\\ c_{t} \CARA (1+\Rfree) & =  - \log (\Rfree\Discount) + \CARA \Rfree m_{t} - (\CARA^{2} \sigma^{2}_{\PShk}/2)
\\ c_{t}  & =  \left(\frac{1}{1+\Rfree}\right)\left[\log (\Rfree\Discount)^{-1/\CARA} + \Rfree m_{t} - \CARA \sigma^{2}_{\PShk}/2 \right]
%\\ \CARA (c_{t+1}-c_{t})  & \approx  r - \tau
%\\ \Delta c_{t+1} & \approx  \CARA^{-1}(\rfree - \tau)
\end{aligned}\end{gathered}\end{equation}

\begin{equation}\begin{gathered}\begin{aligned}
        \uFunc^{\prime}(c_{t}) & =  \Rfree\Discount \Ex_{t}[\uFunc^{\prime}(c_{t+1})] \\
    \exp[-\CARA c_{t}] & =  \Rfree \Discount \Ex_{t}[\exp[-\CARA {c}_{t+1}]]
\\  1 & =  \Rfree \Discount \Ex_{t}[\exp[-\CARA (\left(\frac{1}{1+\Rfree}\right)\left[\log (\Rfree\Discount)^{-1/\CARA} + \Rfree ((m_{t}-c_{t})\Rfree+\PShk_{t+1}) - \CARA \sigma^{2}_{\PShk}/2 \right]-c_{t})]] \nonumber
\\  1 & =  \Rfree\Discount \Ex_{t}\left\{\exp[-\CARA (\left(\frac{1}{1+\Rfree}\right)\left[\log (\Rfree\Discount)^{-1/\CARA} + \Rfree ((m_{t}-c_{t})\Rfree) - \CARA \sigma^{2}_{\PShk}/2 \right]-c_{t})]\exp[\left(\frac{-\CARA \Rfree {\PShk}_{t+1}}{1+\Rfree}\right)] \right\} \nonumber
\\  1 & =  \Rfree\Discount \Ex_{t}\exp[-\CARA (\left(\frac{\log (\Rfree\Discount)^{-1/\CARA} + \Rfree ((m_{t}-c_{t})\Rfree) - \CARA \sigma^{2}_{\PShk}/2}{1+\Rfree}\right)-c_{t})]\exp[\left(\frac{-\CARA \Rfree }{1+\Rfree}\right)^{2}(\sigma_{\PShk}^{2}/2)]  \nonumber
\\  1 & =  \Rfree\Discount \Ex_{t}\exp[-\CARA (\left(\frac{\log (\Rfree\Discount)^{-1/\CARA} + \Rfree^{2} m_{t} - \CARA \sigma^{2}_{\PShk}/2}{1+\Rfree}\right)-c_{t}(1+\frac{\Rfree^{2}}{1+\Rfree})]\exp[\left(\frac{-\CARA \Rfree }{1+\Rfree}\right)^{2}(\sigma_{\PShk}^{2}/2)]  \nonumber
\\  1 & =  \Rfree\Discount \Ex_{t}\exp[-\CARA (\left(\frac{\log (\Rfree\Discount)^{-1/\CARA} + \Rfree^{2} m_{t} - \CARA \sigma^{2}_{\PShk}/2}{1+\Rfree}\right)-c_{t}\left(\frac{1+\Rfree+\Rfree^{2}}{1+\Rfree}\right)]\exp[\left(\frac{-\CARA \Rfree }{1+\Rfree}\right)^{2}(\sigma_{\PShk}^{2}/2)]  \nonumber
\\  0 & =  \log (\Rfree\Discount) + (\CARA c_{t}\left(\frac{1+\Rfree+\Rfree^{2}}{1+\Rfree}\right) - \CARA \left(\frac{\log (\Rfree\Discount)^{-1/\CARA} + \Rfree^{2} m_{t} - \CARA \sigma^{2}_{\PShk}/2}{1+\Rfree}\right) + \left(\frac{-\CARA \Rfree }{1+\Rfree}\right)^{2}(\sigma_{\PShk}^{2}/2) \nonumber
\\  c_{t} & =  \log (\Rfree\Discount)^{-1/\CARA}\left(\frac{1+\Rfree}{1+\Rfree+\Rfree^{2}}\right)  + \CARA \left(\frac{\log (\Rfree\Discount)^{-1/\CARA} + \Rfree^{2} m_{t} - \CARA \sigma^{2}_{\PShk}/2}{1+\Rfree+\Rfree^{2}}\right) - \left(\frac{\CARA \Rfree^{2} (\sigma^{2}_{\PShk})/2}{(1+\Rfree)(1+\Rfree+\Rfree^{2})}\right) \nonumber
\\  c_{t} & =  \log (\Rfree\Discount)^{-1/\CARA}\left(\frac{1+\Rfree}{1+\Rfree+\Rfree^{2}}\right)  + \CARA \left(\frac{\log (\Rfree\Discount)^{-1/\CARA} + \Rfree^{2} m_{t} - \CARA \sigma^{2}_{\PShk}/2}{1+\Rfree+\Rfree^{2}}\right) - \left(\frac{\CARA \Rfree^{2} (\sigma^{2}_{\PShk})/2}{(1+\Rfree)(1+\Rfree+\Rfree^{2})}\right) \nonumber
\\ c_{t} \CARA (1+\Rfree) & =  - \log (\Rfree\Discount) + \CARA \Rfree m_{t} - (\CARA^{2} \sigma^{2}_{\PShk}/2)
\\ c_{t}  & =  \left(\frac{1}{1+\Rfree}\right)\left[\log (\Rfree\Discount)^{-1/\CARA} + \Rfree m_{t} - \CARA \sigma^{2}_{\PShk}/2 \right]
\end{aligned}\end{gathered}\end{equation}

Based on this, let's postulate a consumption rule of the form $c_{t} = 
\mu + \omega m_{t}$ and see whether we can determine coefficients 
$\mu, \omega, $ such that the equation holds in the infinite horizon 
case.

\begin{equation*}\begin{gathered}\begin{aligned}
        0  & =  \log (\Rfree \Discount)+ \log \Ex_{t}\left\{ \exp[-\CARA ({c}_{t+1}-c_{t})]\right\}  \\
        c_{t+1}-c_{t} & =  \mu + \omega((m_{t}-c_{t})\Rfree+y_{t+1})-c_{t}  \\
         & =  \mu + \omega(\Rfree m_{t}+y_{t+1}) - c_{t}(\omega \Rfree + 1)
\\   & =  \mu + \omega \Rfree m_{t} + \omega y_{t+1} - c_{t}(\omega \Rfree + 1)   
\\ \Ex_{t}\left[\exp[-\CARA ({c}_{t+1}-c_{t})]\right] & = \Ex_{t}\left[ \exp[-\CARA (\mu + \omega \Rfree m_{t} + \omega {y}_{t+1} - c_{t}(\omega \Rfree + 1)    ]\right]
\\  & =  \Ex_{t}\left\{ \exp[-\CARA (\mu + \omega \Rfree m_{t} - c_{t}(\omega \Rfree + 1)) -\CARA        \omega {y}_{t+1}  ]\right\}
\\  & =  \exp[-\CARA (\mu + \omega \Rfree m_{t} - c_{t}(\omega \Rfree + 1))]\Ex_{t}\left\{\exp[ -\CARA   \omega {y}_{t+1}  ]\right\}
\\  & =  \exp[-\CARA (\mu + \omega \Rfree m_{t} - c_{t}(\omega \Rfree + 1))]\exp[(\CARA \omega \sigma_{y})^{2}/2]
\\      0  & =  \log (\Rfree \Discount)-\CARA (\mu + \omega \Rfree m_{t} - c_{t}(\omega \Rfree + 1))+(\CARA \omega \sigma_{y})^{2}/2
\\      c_{t}  & =  \left(\frac{-\CARA}{\omega \Rfree + 1}\right)\log (\Rfree \Discount)-\CARA (\mu + \omega \Rfree m_{t} - c_{t}(\omega \Rfree + 1))+(\CARA \omega \sigma_{y})^{2}/2
\end{aligned}\end{gathered}\end{equation*}

\begin{equation*}\begin{gathered}\begin{aligned}
        0  & =  \log (\Rfree \Discount)+ \log \Ex_{t}\left\{ \exp[-\CARA ({c}_{t+1}-c_{t})]\right\}  
\\      c_{t+1}-c_{t} & =  \omega(m_{t+1}-m_{t})
\\       & =  \omega((m_{t}-c_{t})R+y_{t+1}-m_{t})
\\   & =  \omega(\rfree m_{t} - \Rfreec_{t})+\omega y_{t+1}
\\   & =  \omega(\rfree m_{t} - \Rfree(\mu + \omega m_{t})) + \omega y_{t+1}
\\   & =  \omega(m_{t}(\rfree - \Rfree\omega) - \Rfree\mu) + \omega y_{t+1}
\\ c_{t+1} & =  c_{t} + \omega(m_{t}(\rfree - \Rfree\omega) - \Rfree\mu) + \omega y_{t+1}
\\         & =  c_{t} + \omega((c_{t}-\mu)/\omega)(\rfree - \Rfree\omega) - \Rfree\mu) + \omega y_{t+1}
\\         & =  c_{t} + (c_{t}-\mu)(\rfree - \Rfree\omega) - \Rfree\omega\mu + \omega y_{t+1}
\\         & =  c_{t}(1 + (\rfree-\Rfree\omega)) -\m\uFunc(\rfree-\Rfree\omega) -\Rfree\omega\mu + \omega y_{t+1}
\\         & =  c_{t}(1 + (\rfree-\Rfree\omega)) -\mu r + \omega y_{t+1}
\\ c_{t+2} & =  c_{t+1} + \omega(m_{t+1}(\rfree - \omega) - \mu) + \omega y_{t+2}
\\         & =  c_{t}+ \omega(m_{t}(\rfree - \omega) - \mu) + \omega y_{t+1} + \omega(((m_{t}-c_{t})R+y_{t+1})(\rfree - \omega) - \mu) + \omega y_{t+2}
\end{aligned}\end{gathered}\end{equation*}

\begin{equation*}\begin{gathered}\begin{aligned}
        0  & =  \log (\Rfree \Discount)+ \log \Ex_{t}\left\{ \exp[-\CARA ({c}_{t+1}-c_{t})]\right\}  \\
        c_{t+1}-c_{t} & =  \mu + \omega((m_{t}-c_{t})R+y_{t+1})-c_{t}  \\
         & =  \mu + \omega(\Rfree m_{t}+y_{t+1}) - c_{t}(\omega \Rfree + 1)
\\ c_{t} & =    \left(\frac{\mu + \omega( \Rfree m_{t}+y_{t+1})}{\omega \Rfree + 1}\right)
\\       & =    \left(\frac{\mu + \omega \Rfree m_{t}}{\omega \Rfree + 1}\right)+\left(\frac{\omega}{\omega \Rfree + 1}\right)y_{t+1}
\end{aligned}\end{gathered}\end{equation*}


Now for period $t = T-2$ we have
\begin{equation}\begin{gathered}\begin{aligned}
        \uFunc^{\prime}(c_{t}) & =  \Rfree\Discount \Ex_{t}[\uFunc^{\prime}(c_{t+1})] \\
    \exp[-\CARA c_{t}] & =  \Rfree \Discount \Ex_{t}[\exp[-\CARA {c}_{t+1}]]
\\  1 & =  \Rfree \Discount \Ex_{t}[\exp[-\CARA ((m_{t}-c_{t})R+{y}_{t})-c_{t})]]
\\  1 & =  \Rfree\Discount \Ex_{t}[\exp[-\CARA (-c_{t}(1+\Rfree) + Rm_{t})]\exp[-\CARA {y}_{t+1}]]
%\\  1 & =  \log (\Rfree\Discount) + (\CARA c_{t}(1+\Rfree) - \CARA \Rfree m_{t})+ \log \Ex_{t}[\exp[-\CARA {y}_{t+1}]] \nonumber
\\ c_{t} \CARA (1+\Rfree) & =  - \log (\Rfree\Discount) + \CARA \Rfree m_{t} - (\CARA^{2} \sigma^{2}_{y}/2)
\\ c_{t}  & =  \left(\frac{1}{1+\Rfree}\right)\left[\log (\Rfree\Discount)^{-1/\CARA} + \Rfree m_{t} - \CARA \sigma^{2}_{y}/2 \right]
%\\ \CARA (c_{t+1}-c_{t})  & \approx  r - \tau
%\\ \Delta c_{t+1} & \approx  \CARA^{-1}(\rfree - \tau)
\end{aligned}\end{gathered}\end{equation}


\pagebreak Now we need a couple of facts:
\begin{Fact} $\displaystyle \sum_{i=0}^{T} \gamma^{i} = \left(\frac{1-\gamma^{T+1}}{1-\gamma}\right)$ 
\end{Fact}

\begin{Fact} If $0 < \gamma < 1$, then $\displaystyle \sum_{i=0}^{\infty} \gamma^{i} = \left(\frac{1}{1-\gamma}\right)$ \label{fact:fact2}
\end{Fact}

\begin{Fact} If $1 < R$, then $\displaystyle \sum_{i=0}^{\infty} i/\Rfree^{i} = \left(\frac{\Rfree}{(\Rfree-1)^{2}}\right)$ \label{fact:fact3}
\end{Fact}


The Intertemporal Budget Constraint tells us that the present discounted \opt{MarginNotes}{\marginpar{\tiny Note that fact 2 is an obvious implication of fact 1 as $T \rightarrow \infty$.}}{}
value of consumption must be equal to the PDV of total resources:
\begin{equation}\begin{gathered}\begin{aligned}
        \PDV_{t}(\CRat) & =  B_{t}+\PDV_{t}(P)
        \label{eq:ibc}
\end{aligned}\end{gathered}\end{equation}

Using Fact 1, the PDV of labor income (also called `human wealth' $H_{t}$) is
\begin{equation}\begin{gathered}\begin{aligned}
        H_{t} = \PDV_{t}(P) & =  \sum_{s=t}^{T} \Rfree^{-(s-t)}\PLev_{s}
\\       & =  \PLev_{t}\sum_{s=t}^{T} \Rfree^{-(s-t)}\PGro^{(s-t)}
\\       & =  \PLev_{t}\sum_{s=t}^{T} (\PGro/\Rfree)^{(s-t)}
\\   & =  \PLev_{t}\left(\frac{1-(\PGro/\Rfree)^{T-t+1}}{1-(\PGro/\Rfree)}\right)    \label{eq:yfin}
\end{aligned}\end{gathered}\end{equation}
while the PDV of consumption is
\begin{equation}\begin{gathered}\begin{aligned}
        \PDV_{t}(\CRat) & =  \sum_{s=t}^{T} \Rfree^{-(s-t)}\CRat_{s}
\\       & =  \CRat_{t}+(\CRat_{t}+ \log (\Rfree\Discount)^{-1/\CARA})/\Rfree +(\CRat_{t}+ 2 \log (\Rfree\Discount)^{-1/\CARA})/\Rfree^{2} + \ldots
\\       & =  \CRat_{t}(1+1/\Rfree+1/\Rfree^{2}+\ldots) + \log (\Rfree\Discount)^{-1/\CARA}/\Rfree + 2 \log (\Rfree\Discount)^{-1/\CARA}/\Rfree^{2} + \ldots \nonumber
\\       & =  \left(\frac{\CRat_{t}}{1-1/\Rfree}\right) +  \log (\Rfree\Discount)^{-1/\CARA}(1/\Rfree + 2/\Rfree^{2} + 3/\Rfree^{3} + \ldots)
\\   & =  \left(\frac{\CRat_{t}}{1-1/\Rfree}\right) +  \log (\Rfree\Discount)^{-1/\CARA}\left(\frac{\Rfree}{(\Rfree-1)^{2}}\right)
\end{aligned}\end{gathered}\end{equation}

Therefore we can solve the model by combining \eqref{eq:cfin} and \eqref{eq:yfin} using 
\eqref{eq:ibc}:
\begin{equation*}\begin{gathered}\begin{aligned}
        \left(\frac{\CRat_{t}}{1-1/\Rfree}\right) +  \log (\Rfree\Discount)^{-1/\CARA}\left(\frac{\Rfree}{(\Rfree-1)^{2}}\right) & =  B_{t}+\left(\frac{\PLev_{t}}{1-\PGro/\Rfree}\right)  \\
        \left(\frac{\CRat_{t}}{1-1/\Rfree}\right) & =  B_{t}+\left(\frac{\PLev_{t}}{1-\PGro/\Rfree}\right) -\log (\Rfree\Discount)^{-1/\CARA}\left(\frac{\Rfree}{(\Rfree-1)^{2}}\right) \\     
        \CRat_{t} & =  \left(\frac{\Rfree}{\rfree}\right)\left[B_{t}+\left(\frac{\PLev_{t}}{1-\PGro/\Rfree}\right) -\log (\Rfree\Discount)^{-1/\CARA}\left(\frac{\Rfree}{(\Rfree-1)^{2}}\right)\right]
\end{aligned}\end{gathered}\end{equation*}
\begin{equation}\begin{gathered}\begin{aligned}
        \CRat_{t} & =  \left(\frac{1-[\Rfree^{-1}(\Rfree\Discount)^{1/\CARA}]}{1-[\Rfree^{-1}(\Rfree\Discount)^{1/\CARA}]^{T+1}}\right)\left[B_{t}+\PLev_{t}\left(\frac{1-(\PGro/\Rfree)^{T+1}}{1-(\PGro/\Rfree)}\right)\right] \label{eq:cfinhoriz}
\end{aligned}\end{gathered}\end{equation}

Now recall that in the infinite-horizon case ($T=\infty$), Fact 2 
requires that for human wealth to be well-defined we need the 
condition
\begin{equation}\begin{gathered}\begin{aligned}
        \PGro/\Rfree & <  1
\\  \PGro   & <  R.  \label{eq:pdvyinf}
\end{aligned}\end{gathered}\end{equation}
Why is this?  Because if income will grow faster than the interest 
rate forever, then the PDV of future income is infinite and the 
problem has no well-defined solution.

Similarly, in order for the PDV of consumption to be finite we must
impose:
\begin{equation}\begin{gathered}\begin{aligned}
        \Rfree^{-1}(\Rfree\Discount)^{1/\CARA}& < 1
\\      (\Rfree\Discount)^{1/\CARA} & <  R. \label{eq:pdvcinf}
\end{aligned}\end{gathered}\end{equation}
What this says is that the growth rate of consumption must be less 
than the interest rate in order for the model to have a well-defined 
solution.  Otherwise, the PDV of future consumption is infinite, and 
the model does not have a well-defined solution.  Note that this 
amounts to a requirement that there be at least a certain degree of 
`impatience.'

If these conditions do hold, then the model has a well-defined 
infinite horizon solution, as can be seen by realizing that if 
$(\PGro/\Rfree)<1$ then $\lim_{T \rightarrow \infty} (\PGro/\Rfree)^{T-t+1} = 0$ and if 
$\Rfree^{-1}(\Rfree\Discount)^{1/\CARA} < 1$ then $\displaystyle \lim_{T \rightarrow 
\infty} (\Rfree^{-1}(\Rfree\Discount)^{1/\CARA})^{T-t+1} = 0$.  Substituting these 
zeros into \eqref{eq:cfinhoriz} yields 
\begin{equation}\begin{gathered}\begin{aligned}
        \CRat_{t} & =  \left(1-\Rfree^{-1}(\Rfree\Discount)^{1/\CARA}\right)\left[B_{t}+\left(\frac{\PLev_{t}}{1-(\PGro/\Rfree)}\right)\right] \label{eq:cinfhor}
\\      & =  \left(1-\Rfree^{-1}(\Rfree\Discount)^{1/\CARA}\right)\left[B_{t}+H_{t}\right]  \label{eq:cmpk}
\\      & =  \left(\frac{\Rfree-(\Rfree\Discount)^{1/\CARA}}{\Rfree}\right)W_{t} \label{eq:cOfw}
\end{aligned}\end{gathered}\end{equation}
where $W_{t}$ is the consumer's `total wealth,' the sum of human and 
nonhuman wealth.

Now consider the question `What is the level of $\CRat_{t}$ that will 
leave total wealth intact, allowing the same value of consumption in 
period $t+1$ and forever after?'

The intuitive answer is that if one wants to leave one's wealth 
\opt{MarginNotes}{\marginpar{\tiny Note that this was 
interpreted as `permanent income' in the 1960s and 70s, but will not 
be called such in this class.  Point out that wealth here is exactly 
like an asset that yields a dividend $P$.}}{} intact, that is possible 
only if spending is exactly equal to the dividend and interest 
earnings on one's total wealth.

Because human wealth is exactly like any other kind of wealth in this 
framework, it is possible to work directly with the level of total 
wealth $W$.  Suppose we assume the consumer will spend fraction 
$\kappa$ of total wealth in each period, and we want to find the 
$\kappa$ that leaves wealth intact.
\begin{equation}\begin{gathered}\begin{aligned}
        W_{t+1} & =  (W_{t}-\CRat_{t})R
\\  \bar{W} & =  (\bar{W}-\kappa \bar{W})R     
\\      1 & =  \Rfree(1-\kappa)  
\\      1/\Rfree & =  (1-\kappa)
\\  \kappa & =  1-1/\Rfree
\\  & =  \left(\frac{\Rfree-1}{\Rfree}\right)
\\  & =  \rfree/\Rfree \label{eq:kappa}
\end{aligned}\end{gathered}\end{equation}
Thus, the consumer can spend only the interest earnings $\rfree$ on their 
wealth, divided by the gross return $\Rfree$.  (The division occurs because 
we assume that interest is earned between periods rather than within 
periods; the right intuition is that if you want to preserve your 
wealth, you can only spend the interest on it and none of the 
principal).


Note that the coefficient multiplying total wealth in \eqref{eq:cOfw} 
is also divided by $\Rfree$.  Thus, whether the consumer is spending more 
than his total income, exactly his total income, or less than his 
total income depends upon whether the numerator in \eqref{eq:cOfw} is
greater than, equal to, or less than $\rfree$.  If we call a consumer 
who is spending more than his income `impatient,' the consumer will be
impatient if
\begin{equation}\begin{gathered}\begin{aligned}
        R-(\Rfree\Discount)^{1/\CARA} & >  r  \\
        1-(\Rfree\Discount)^{1/\CARA} & >  0  \\
        1 & >  (\Rfree\Discount)^{1/\CARA}
\end{aligned}\end{gathered}\end{equation}

Now note that if $\Rfree\Discount=1$ (which is to say, the interest rate is 
exactly equal to the time preference rate so that they offset each 
other), then $(\Rfree\Discount)^{1/\CARA}=1$ regardless of the value of $\CARA$ 
so that the consumer is precisely poised on the balance between 
patience and impatience and exactly spends his income.\opt{MarginNotes}{\marginpar{\tiny Income here means inclusive of interest income on total wealth.}}{}

The consumer will be impatient, spending more than his income, if 
$\Rfree\Discount>1$, and patient, spending less than his income, if $\Rfree\Discount<1$.

Equation \eqref{eq:cinfhor} can be simplified into something a bit 
easier to handle by making some approximations.  If $\Discount = 
1/(1+\tau)$, then we can use 
\begin{Fact} $\log (1+\PShk) \approx 
\PShk$
\end{Fact}
and its inverse
\begin{Fact} $\exp(\PShk) \approx 1+\PShk$
\end{Fact}
to discover that
\begin{equation}\begin{gathered}\begin{aligned}
        \log (\Rfree\Discount)^{1/\CARA}/\Rfree & =  (1/\CARA) (\log \Rfree + \log [1/(1+\tau) ]) - \log \Rfree  \\
     & =  (1/\CARA) (\log(1+r) + \log 1 - \log (1+\tau) ) - \log \Rfree  \\
         & \approx  \CARA^{-1}(\rfree -\tau) ) - r 
\\      (\Rfree\Discount)^{1/\CARA}/\Rfree & \approx  1+(\CARA^{-1}(\rfree-\tau)-\rfree)
\end{aligned}\end{gathered}\end{equation}

Substituting this into \eqref{eq:cmpk} gives
\begin{equation}\begin{gathered}\begin{aligned}
        \CRat_{t} & \approx  \left(\rfree-\CARA^{-1}(\rfree-\tau)\right)W_{t} \label{eq:capprox}
\end{aligned}\end{gathered}\end{equation}

Now we can see again that whether the consumer is patient or impatient 
depends on the relationship between $\rfree$ and $\tau$.  Note also that 
if $\CARA = \infty$ then the consumer is infinitely averse to changing 
the level of consumption, and so once again the consumer spends 
exactly his income.


Now a  brief note on what `income' means in this model.  Suppose
for simplicity that the consumer had no capital assets $K$, and suppose
that income was expected to stay constant at level $\bar{Y}$ forever.
In this case human wealth would be:
\begin{equation}\begin{gathered}\begin{aligned}
        H_{t} & =  \bar{Y}+\bar{Y}/\Rfree+\bar{Y}/\Rfree^{2}+\ldots  \\
         & =  \bar{Y}(1+1/\Rfree+1/\Rfree^{2}+\ldots)  \\
         & =  \bar{Y}\left(\frac{1}{1-1/\Rfree}\right)
\\   & =  \bar{Y}\left(\frac{\Rfree}{\Rfree-1}\right)
\\   & =  \bar{Y}\left(\frac{\Rfree}{\rfree}\right)
\end{aligned}\end{gathered}\end{equation}

Now recall that we found in equation \eqref{eq:kappa} that 
the level of consumption that leaves `wealth' $W_{t}$ intact
was
\begin{equation}\begin{gathered}\begin{aligned}
        \CRat_{t} & =  \kappa W_{t}  \\
         & =  \kappa [B_{t}+H_{t}] \\
         & =  \kappa \bar{Y}\left(\frac{\Rfree}{\rfree}\right)
\\   & =  \left(\frac{\rfree}{\Rfree}\right)     \bar{Y} \left(\frac{\Rfree}{\rfree}\right)
\\   & =  \bar{Y}.
\end{aligned}\end{gathered}\end{equation}

So in this case, spending the `interest income on human wealth' 
corresponds to spending exactly your labor income.  This seems less 
mysterious if you think of income $Y_{t}$ as the `return' on your 
human capital asset $H_{t}$.  If you `capitalize' your stream of 
income at rate $\Rfree$ and then spend the interest income on the 
capitalized stream, it stands to reason that you are spending the flow 
of income from that source.

Note also that in this case we can rewrite \eqref{eq:capprox} as
\begin{equation}\begin{gathered}\begin{aligned}
        \CRat_{t} & \approx  \left(\rfree-\CARA^{-1}(\rfree-\tau)\right)\left[B_{t}+ \bar{Y}\left(\frac{\Rfree}{\rfree}\right)\right].
\end{aligned}\end{gathered}\end{equation}

Note that $\rfree$ appears three times in this equation, which correspond 
(in order) to the income effect, the substitution effect, and the 
human wealth effect.  To see this, note that an increase in the first 
$\rfree$ basically corresponds to an increase in the payout rate on total 
wealth (to see this, set $\bar{Y} = 0$ and refer to our formula above 
for $\kappa$, realizing that for small $\rfree$, $\rfree/\Rfree \approx \rfree$.)  The 
second term corresponds to the subsitution effect, as can be seen from 
its dependence on the intertemporal elasticity of substitution $\CARA^{-1}$.  
Finally, the $\bar{Y}/\rfree$ term clearly corresponds to human wealth, and 
therefore the sensitivity of consumption to $\rfree$ coming through this 
term corresponds to the human wealth effect.  \opt{MarginNotes}{\marginpar{\tiny This is why I introduced the concept of the human wealth effect in my original treatment in the Fisher diagram.}}{}



\vspace{.3in} \noindent {\bf Renormalizing by $P$}

Finally, note that we can restate the whole problem by `dividing 
through' by the level of permanent income before solving.  That is, 
define small-letter variables as the capital-letter equivalent, 
divided by the level of permanent labor income in the corresponding 
period, e.g.  $c_{t}=\CRat_{t}/\PLev_{t}$, and note that if $\PLev_{s+1}=\PGro 
\PLev_{s}~\forall~s$ then from the standpoint of time $t$ we have that
\begin{equation}\begin{gathered}\begin{aligned}
        \uFunc(\CRat_{s}) & =  \frac{\CRat_{s}^{1-\CARA}}{1-\CARA}  \\
         & =  \frac{(c_{s}\PLev_{s})^{1-\CARA}}{1-\CARA}  \\
         & =  (\PLev_{t}\PGro^{s-t})^{1-\CARA}\frac{c_{s}^{1-\CARA}}{1-\CARA}
\end{aligned}\end{gathered}\end{equation}
which means that 
\begin{equation}\begin{gathered}\begin{aligned}
        \sum_{s=t}^{T} \Discount^{(s-t)}\frac{\CRat_{s}^{1-\CARA}}{1-\CARA} & =  \PLev_{t}^{1-\CARA}\sum_{s=t}^{T} (\PGro^{1-\CARA}\Discount)^{(s-t)} \frac{c_{s}^{1-\CARA}}{1-\CARA}  \label{eq:maxc2}
\end{aligned}\end{gathered}\end{equation}

Furthermore, the accumulation equations can be rewritten by dividing
both sides by $\PLev_{t+1}$:
\begin{equation}\begin{gathered}\begin{aligned}
        B_{t+1}/\PLev_{t+1} & =  \frac{({M}_{t}-\CRat_{t})\Rfree}{\PLev_{t+1}}  \\
        k_{t+1} & =   \left(\frac{({M}_{t}-\CRat_{t})\Rfree}{\PLev_{t}}\right)\left(\frac{\PLev_{t}}{\PLev_{t+1}}\right) \\
         & =  (\Rfree/\PGro)[m_{t}-c_{t}]
\end{aligned}\end{gathered}\end{equation}
\begin{equation}\begin{gathered}\begin{aligned}
        {M}_{t+1} & =  B_{t+1}+\PLev_{t+1}
\\      m_{t+1} & =  k_{t+1}+1
\end{aligned}\end{gathered}\end{equation}

Now if we define $\hat{\Discount} = \PGro^{1-\CARA}\Discount$ and $\hat{\Rfree} = \Rfree/\PGro$, and 
$\hat{\PGro}=1$, it is clear that the original problem can be rewritten as:


\begin{equation}
\max \PLev_{t}\sum_{s=t}^{T} \hat{\Discount}^{s-t} \uFunc(c_{s})             \label{eq:scaledmaxprob}
\end{equation}
subject to the constraints
\begin{equation}\begin{gathered}\begin{aligned}
    k_{t+1} & =  \hat{\Rfree}[m_{t}-c_{t}]
\\      m_{t+1} & =  k_{t+1}+1
\end{aligned}\end{gathered}\end{equation}
and we can go through the same steps as above to find that the solution
is
\begin{equation}\begin{gathered}\begin{aligned}
        c_{t} & =  (1-\hat{\Rfree}^{-1}(\hat{\Rfree}\hat{\Discount})^{1/\CARA})\left[k_{t}+\left(\frac{1}{1-\hat{\PGro}/\hat{\Rfree}}\right)\right] \label{eq:lowerc}
\end{aligned}\end{gathered}\end{equation}
subject to the `finite human wealth' condition
\begin{equation}\begin{gathered}\begin{aligned}
        \hat{\PGro} & <  \hat{\Rfree}
\\  1 & <  \Rfree/\PGro 
\end{aligned}\end{gathered}\end{equation}
which is the same condition as above \eqref{eq:pdvyinf}, and the `impatience condition'
\begin{equation}\begin{gathered}\begin{aligned}
                (\hat{\Rfree}\hat{\Discount})^{1/\CARA} & <  \hat{\Rfree}  
\\              \left(\frac{\Rfree}{\PGro}\Discount \PGro^{1-\CARA}\right)^{1/\CARA} & <  \Rfree/\PGro  
\\              (\Rfree\Discount)^{1/\CARA} & <  \Rfree
\end{aligned}\end{gathered}\end{equation}  
which is the same as \eqref{eq:pdvcinf}.

Now note that \eqref{eq:lowerc} can be rewritten
\begin{equation}\begin{gathered}\begin{aligned}
        c_{t} & =  \left(\frac{\hat{\Rfree}-(\hat{\Rfree}\hat{\Discount})^{1/\CARA}}{\hat{\Rfree}}\right)w_{t}
\end{aligned}\end{gathered}\end{equation}
where $w_{t}$ is the consumer's total wealth-to-permanent-labor-income 
ratio, human and nonhuman.

As before, whether $w$ is rising or falling depends upon the 
relationship between $\hat{\rfree}$ and 
$\hat{\Rfree}-(\hat{\Rfree}\hat{\Discount})^{1/\CARA}$.  If we call a consumer who is 
drawing down his wealth-to-income ratio `impatient,' the consumer will 
be impatient if
\begin{equation}\begin{gathered}\begin{aligned}
        R-(\Rfree\Discount)^{1/\CARA} & >  r  \\
        1-(\Rfree\Discount)^{1/\CARA} & >  0  \\
        1 & >  (\hat{\Rfree}\hat{\Discount})^{1/\CARA}
\end{aligned}\end{gathered}\end{equation}

Now substituting the definitions of $\hat{\Rfree}$ and $\hat{\Discount}$
we see that whether $w$ is rising or falling depends on whether
\begin{equation}\begin{gathered}\begin{aligned}
        1 & >  (\frac{\Rfree}{\PGro}\Discount \PGro^{1-\CARA})^{1/\CARA}
\\  1 & >  (\Rfree\Discount \PGro^{-\CARA})^{1/\CARA}
\\  1 & >  \PGro^{-1}(\Rfree\Discount)^{1/\CARA}
\\  \PGro & >  (\Rfree\Discount)^{1/\CARA}
\end{aligned}\end{gathered}\end{equation}

Thus, whether the consumer is patient or impatient in the sense of 
building up or drawing down a wealth-to-income ratio depends on 
whether the growth rate of labor income is less than, equal to, or 
greater than the growth rate of consumption.  

To get the intuition for this, consider the case of a consumer with no 
nonhuman wealth, $k_{t}=0$.  This consumer's absolute level of 
consumption will grow at $(\Rfree\Discount)^{1/\CARA}$ and absolute level of 
income at $\PGro$, but the PDV of future consumption and future income are 
equal.  Thus, if income is growing faster than consumption but has the 
same PDV, consumption must be starting out at a level higher than 
income - which is to say, if the impatience condition holds, the 
consumer has a high level of consumption but slow consumption growth.

A final point about the human wealth effect.  For simplicity, assume 
that $k_{t} = 0$.  Then the original version of the formula informs
us that the {\it level} of consumption will be given by:
\begin{equation}\begin{gathered}\begin{aligned}
        \CRat_{t} & \approx  \left[\rfree - \CARA^{-1}(\rfree-\tau)\right]\left(\frac{\bar{Y}}{1-\Rfree/\PGro}\right)  \\
         & \approx  \left[\rfree - \CARA^{-1}(\rfree-\tau)\right]\left(\frac{\bar{Y}}{\rfree-g}\right). \label{eq:cofw2}
\end{aligned}\end{gathered}\end{equation}

Now suppose we choose plausible values for $(\rfree, \tau, g, \CARA) = (.04,.04,.02,2)$.
Then \eqref{eq:cofw2} becomes:
\begin{equation}\begin{gathered}\begin{aligned}
        \CRat_{t} & \approx  .04 (\bar{Y}/.02) \\
         & =  2 \bar{Y}.
\end{aligned}\end{gathered}\end{equation}

Now suppose the interest rate changes to $\rfree=.03$, while all other parameters
remain the same.  Then \eqref{eq:cofw2} becomes:
\begin{equation}\begin{gathered}\begin{aligned}
        \CRat_{t} & \approx  .025 (\bar{Y}/.01)  \\
         & =  2.5 \bar{Y}.
\end{aligned}\end{gathered}\end{equation}

The point of this analysis is that for plausible parameter values, the 
human wealth effect is enormously stronger than the income and 
substitution effects, so that we should see large drops in consumption 
when interest rates rise and conversely strong gains when interest 
rates fall.  This is a summary of the main point of the paper by 
Summers~\cite{summersCapTax} on your syllabus; Summers derives 
formulas for an economy with overlapping generations of 
finite-lifetime consumers, but those complications do not change the 
basic message.

%\input texhtml.tex
\bibliographystyle{econometrica}
\input handoutBibMake


\end{document}

  \\
         & \mbox{s.t.}  \nonumber \\
        {M}_{t+1} & =  ({M}_{t}-\CRat_{t})R+Y_{t+1}
        \label{eq:budconstr}
\end{aligned}\end{gathered}\end{equation}
